\clearpage
\begingroup
\fancyhead[L]{\small POSTHUMOUS PAPERS AND CORRESPONDENCE}
\renewcommand{\headrulewidth}{0pt}
\par\medskip
\Needspace{8\baselineskip}
\pdfbookmark[2]{[5] Nicolai to Gauss, 6 August 1814}{pallas-VIII-letter5}
\begin{center}
[5.]\par\medskip
\textsc{Nicolai} to \textsc{Gauss}, Seeberg, 6 August 1814.
\end{center}
\par\noindent\hspace*{\parindent}\{.\ .\ .\ .\ .\quad As for the perturbations of Pallas by Saturn, I have
yet to check them by calculating a few special perturbations
(since my return here I have done hardly any
numerical calculations, having occupied myself instead with some theoretical matters),
but I have gone over the whole course of the calculation once more
with care. I found no errors in it, as I might well have
predicted, knowing the care with which I had carried out
the whole calculation. The assumed sign
of the forces is also quite correct, for afterwards,
I no longer remember why, I too started from the assumption
\(\dfrac{x'-x}{\rho^3}-\dfrac{x'}{r^{\prime3}}\). I hope that for all the elements
you will find your coefficients in agreement with mine. I take further
reassurance as to the correctness of my equations from the fact
that the larger ones agree fairly well with my earlier calculation, which
\SourceLeaf{591}
was, after all, carried out by an entirely different method.
Still, I shall shortly calculate a couple of special perturbations from one opposition
to the next, and I hope to obtain close agreement
here too. Should this fail to occur, I would still be more inclined
to lay the blame on the numerical calculation of the equations than on
the equations themselves. Who knows whether the perturbations due to Mars might also worsen
the agreement of the Pallas oppositions observed so far, if only
slightly?\quad .\ .\ .\ .\ .\}

\par\medskip
\Needspace{8\baselineskip}
\pdfbookmark[2]{[6] Nicolai to Gauss, 18 March and 17 July 1815}{pallas-VIII-letters6}
\begin{center}
[6.]\par\medskip
\textsc{Nicolai} to \textsc{Gauss}, Seeberg, 18 March 1815.
\end{center}
\par\noindent\hspace*{\parindent}\{.\ .\ .\ .\ .\quad I shall shortly carry out your request to combine the two parts
of the perturbation of the epoch for Pallas and send you
the formulas. For the perturbation due to Jupiter this must have been quite
a neat piece of work.\quad .\ .\ .\ .\ .\}

\par\medskip
\Needspace{6\baselineskip}
\begin{center}
\textsc{Nicolai} to \textsc{Gauss}, Seeberg, 17 July 1815.
\end{center}
\par\noindent\hspace*{\parindent}\{Yesterday afternoon I worked on combining the two
parts of the perturbation of the epoch of Pallas by Saturn, and I
hasten to send you the equations I found. The equations
in the first part of the epoch that had no corresponding equations in the second
were all below \(0\overset{\prime\prime}{,}1\), so I have left them out of the
final equations for the epoch. In this new form
of the perturbation of the epoch, however, many a small equation \textls[100]{above} \(0\overset{\prime\prime}{,}1\)
whose separate parts are nonetheless \textls[100]{below} \(0\overset{\prime\prime}{,}1\) will be missing. This will often be the case
especially in your perturbations due to Jupiter. --- In the course of this work
I noticed that if the first part of the perturbation of the epoch is brought
into the same form as the second by adding \(90^0\) to the constant angle,
namely \(a\sin(\varphi+A)\), then, especially in the larger equations,
those constant angles differ little from one another. Is this
phenomenon mere chance, or does it have some theoretical basis?
--- and does something similar occur in the perturbations due to Jupiter?\quad .\ .\ .\ .\ .

\SourceLeaf{592}
\par\medskip
\Needspace{13\baselineskip}
\begin{center}
Complete perturbation of the epoch of Pallas by Saturn.\par
Annual change \(=-1\overset{\prime\prime}{,}31973\)
\end{center}
\begingroup
\small
\newcommand{\VIIIpallas}{\text{\AstroSymbols ⚴}}
\newcommand{\VIIIsaturn}{\text{\AstroSymbols ♄}}
\setlength{\arraycolsep}{2pt}
\renewcommand{\arraystretch}{1.25}
\[
\begin{array}{r@{\,}l@{\qquad}r@{\,}l}
0\overset{\prime\prime}{,}83&\sin(\VIIIpallas+175^0\ 20')&+0\overset{\prime\prime}{,}50&\sin(3\VIIIpallas-2\VIIIsaturn+63^0\ 14')\\
+0{,}53&\sin(2\VIIIpallas+182\ 41)&+0{,}30&\sin(3\VIIIsaturn+207\ 41)\\
+1{,}78&\sin(\VIIIsaturn+140\ 25)&+2{,}41&\sin(\VIIIpallas-3\VIIIsaturn+169\ 38)\\
+0{,}47&\sin(\VIIIpallas+\VIIIsaturn+180\ 3)&+0{,}71&\sin(2\VIIIpallas-3\VIIIsaturn+162\ 17)\\
+0{,}61&\sin(\VIIIpallas-\VIIIsaturn+179\ 37)&+0{,}66&\sin(3\VIIIpallas-3\VIIIsaturn+172\ 34)\\
+0{,}18&\sin(2\VIIIpallas-\VIIIsaturn+161\ 42)&+0{,}46&\sin(\VIIIpallas-4\VIIIsaturn+187\ 14)\\
+2{,}01&\sin(2\VIIIsaturn+221\ 24)&+0{,}24&\sin(2\VIIIpallas-4\VIIIsaturn+164\ 16)\\
+7{,}26&\sin(\VIIIpallas-2\VIIIsaturn+124\ 36)&+0{,}15&\sin(3\VIIIpallas-4\VIIIsaturn+161\ 42).\}\\
+3{,}69&\sin(2\VIIIpallas-2\VIIIsaturn+178\ 30)&&
\end{array}
\]
\endgroup
\begin{center}
\rule{0.3\linewidth}{0.4pt}
\end{center}
\clearpage
\endgroup
